\subsection{The two-stage complex endpoint correction}
\label{sec:two-stage-phase-transfer}
Use the original phase operators $C_U=H^{\otimes m}\operatorname{diag}
(i^{\mathrm{wt}(p_U x)})H^{\otimes m}$. Orthogonal weight additivity
modulo four gives $C_{U\perp V}=C_UC_V$. For a two-triple indicator
$w$, $\mathrm{wt}(w)=9$, so on each block
$P=C_{\langle w\rangle}=C_w=aI+bX_w$, where
$a=(1+i)/2$, $b=(1-i)/2$, and $P^2=X_w$, $P^4=I$.
Let $F=C_{\mathcal F}$ and $E=C_{U_a^\perp}=FP^{-1}$.
The common-frame identity for the two scalar stages gives
\[
 A=-Fy,\qquad B=EP^{-1}x+Ey=FP^{-2}x+Ey.
\]
Copy $A$, apply $P^{-1}$ to the copy, and add it into $B$. Since
$P^{-1}F=E$, the $y$ term cancels. The retained binary basis-change
construction implements $P^{-1}$ with one inverse recursive phase child,
with the same selected axes, spectators and compact reservations. Its
inverse wrapper is the retained diagonal-sign and fourth-root identity.
Copy, scalar addition and cleanup are charged at this node; the unmodified
stream $A$ is parked during the child, as in the direct-swap copy proof.

It remains to remove $P^{-2}$ without assuming a free translation. Write
$Z_w$ for the diagonal sign $(-1)^{w\cdot x}$. On one coordinate,
$ZCZ=iC^{-1}$, and $C^{-1}=CX$. Therefore, with $k=\mathrm{wt}(w)$,
\[
 Z_wFZ_w=i^kFX_w,\qquad
 Z_wX_w=(-1)^kX_wZ_w,\qquad
 i^k Z_wFP^{-2}Z_w=F.
\]
For $f$ selected blocks the same identity has scalar $i^{9f}$ and the
product of the block signs. Apply that diagonal sign to the physical $x$
input before the network, and to its corrected $B$ output afterward,
together with $i^{9f}$. Negate $A$ and undo the bank exchange in the row
merge. Both data roles now receive exactly $F$. Every auxiliary receives
$F$ already, since it is restored logically and has zero/full terminal
labels. These diagonal multipliers are computed sequentially from the
selected address bits; their descriptor work is charged in the retained
record regime. They do not permute addresses.

All ordinary edges are implemented by the original signed rank-one phase
kernel construction for their orthonormal residual bases. The extra $N$
copies use exactly $N$ additional phase children, already included in $s$.
The original phase recursion's proof now uses this completed array identity
in place of its three-stage weight-27 endpoint identity. Its basis-change,
row split, dirty-control reservation and depth-first fixed-tape scheduling
arguments are otherwise the retained ones. The explicit coefficient-node
allowance also includes copies, adds, input/output signs, fourth roots and
inverse wrappers, so the generalized stopped-depth guard uses the complete
charged $s$, not the optimistic edge-only count.
